\documentclass[12pt]{book}
\usepackage[utf8]{inputenc}
\usepackage[T1]{fontenc}
\usepackage{amsmath,amssymb,amsfonts}

\begin{document}

\section{ERRATA}

Page
10, line I: read "see Appendix" instead of "unpublished".

Page
85, Proposition 1.1: For the case of $\mathrm{Qco}(X)$, one needs the hypothesis "$X$" quasi-separated".

Page
88, line 5: read "proper" instead of "of finite type".

Page
120, Ch II §7: It has recently come to my attention that S. Kleiman has independently arrived at the results of [II 7.8] and [II 7.11] (unpublished).

Page
137, line 5: refer to [EGIV 16].

Page
144, line 4: read "invertible sheaf shifted $n$: $n=\mathrm{rel.dim}\,X/Y$" instead of places to the "invertible sheaf".

Page
185, line 9: observe that the condition "$f$ finite" is a consequence of the other conditions.

Page
190, line 7: read "a unique isomorphism" instead of "an isomorphism".

Page
199: add at bottom: "Furthermore, the residue symbol is uniquely characterized by the properties (R0), (R1), (R2), (R5), (R6), and (R7)."

Page
249, line 8: "is" instead of "in".

Page
254, line 8: "77." instead of "§".

Page
288, Theorem 8.3: One must assume also that the fibres of $f$ are of bounded dimension, so that $f^{*}(R^{*})$ will be in $D^{+}$.

Page
297, line 7: read "strengthens " instead of "generalizes".

Page
298, line 7: insert "shifted" after "invertible sheaf".

Page 301, line 5: "not" instead of "now".

Page 301, Problem 2: Giraud [6] provides us with an element of
\[H^{2}\left(X, \mathbb{G}_{m}\right)\]
(cohomology in the Zariski topology), whose vanishing is a necessary and sufficient condition for the global existence of a dualizing complex, supposed to exist locally. On the other hand, it is easy enough to construct a noetherian scheme of finite Krull dimension, which has a dualizing complex locally, but none globally, due to the lack of a global codimension function.

Page 373, line 1: change reference to [EGA V].

Page 403, last line: read "Publ. Math. I.H.E.S. no 29".

\end{document}